\section{设计背景}

\subsection{选题背景}

随着数字内容产业的快速发展，
三维模型在游戏开发、影视制作、虚拟现实、增强现实以及数字孪生等领域中的应用日益广泛。\cite{赵沁平2009虚拟现实综述}
在这些应用场景中，大量三维资产需要在较短时间内完成建模与材质制作，以满足不断增长的内容生产需求。
然而，传统三维资产制作流程通常依赖专业建模软件与人工绘制纹理，
这一过程不仅耗时较长，而且对设计人员的技术经验要求较高，
在复杂项目中往往成为影响生产效率的重要因素。

在传统建模流程中，三维模型通常需要经历建模、UV 展开、纹理绘制与材质调试等多个步骤。
其中，纹理与材质制作往往需要设计人员在专业软件中进行手工绘制或贴图调整，
不仅操作流程复杂，而且在大型项目中需要投入大量人力成本。\cite{陈娜娜2025}
随着三维资产需求的持续增长，
如何利用自动化技术提高三维内容生产效率，
逐渐成为计算机图形学与人工智能领域的重要研究方向。

近年来，深度学习与生成式人工智能技术的发展为三维内容生成提供了新的解决思路。
特别是扩散模型与多模态生成模型的出现，
使得从文本或图像自动生成三维模型及其纹理成为可能。
一些研究工作通过结合可微渲染、扩散模型以及多视角一致性优化等技术，
实现了从二维信息到三维纹理的自动化生成，
显著降低了三维建模的技术门槛。

与此同时，越来越多的 AI 三维生成平台开始以 Web 形式提供在线服务，
用户可以通过浏览器直接完成模型生成与结果预览，
无需安装复杂的软件环境。
这种轻量化的交互方式进一步推动了三维内容生产工具的普及，
使非专业用户也能够参与到三维内容创作过程中。\cite{刘钊2023}

然而，在实际应用中仍然存在一个重要问题：
大量已有三维模型仅包含几何结构而缺乏材质信息，
即所谓的“白模”。
在游戏开发、虚拟场景构建以及数字内容制作等项目中，
设计人员往往需要在同一场景中使用大量不同来源的三维模型。
为了保证整体视觉效果的一致性，
这些模型通常需要采用统一或相近的材质风格。

在这种情况下，
开发人员更希望能够将已有模型的材质风格迁移到其他白模上，
从而实现模型之间的材质统一，
而不是重新生成新的三维模型。
通过这种“模型到模型”的材质迁移方式，
可以在保留原有几何结构的基础上快速调整视觉风格，
从而满足游戏场景、虚拟世界或数字资产库中对风格一致性的需求。

然而，当前多数三维生成平台仍主要关注从文本或图像生成完整三维模型的能力，
而针对已有模型之间的材质迁移功能仍然相对有限。
因此，研究一种能够利用参考模型信息为白模自动生成或迁移材质的系统，
对于提升三维资产生产效率、实现大规模三维内容风格统一具有重要的研究价值与应用意义。

\subsection{设计目标和意义}

本课题旨在设计并实现一种基于参考模型的三维材质自动生成与迁移系统，
通过分析已有带材质模型的纹理信息与几何结构，
为目标白模生成与其结构相匹配的材质贴图，
并构建一个基于 Web 的交互平台，
为用户提供模型上传、材质生成与三维预览等功能，
从而形成完整的自动化三维材质生成流程。

从行业应用角度来看，
随着游戏产业、影视制作、虚拟现实以及元宇宙等领域的发展，
三维资产的需求量呈现快速增长趋势。
在大型项目中，
设计团队往往需要制作大量具有不同风格和材质的三维模型。
如果每个模型都依赖人工绘制纹理，
将会消耗大量时间与人力成本。
通过引入自动化材质生成与迁移技术，
可以在已有模型基础上快速生成符合结构特征的纹理，
从而显著提升三维资产生产效率，
为数字内容产业提供更加高效的生产工具。

从科研角度来看，
三维纹理生成与材质迁移是计算机图形学与人工智能交叉领域的重要研究方向之一。\cite{chen2024advances3dneuralstylization}
现有研究多集中在从文本或图像生成三维模型，
而针对已有几何结构的纹理生成问题仍有较大的研究空间。
通过研究参考模型驱动的材质迁移方法，
可以进一步探索几何结构与纹理分布之间的关系，
并为三维生成模型的研究提供新的思路。
同时，该研究也有助于推动可微渲染、生成模型与三维表示学习等技术的进一步融合。

从技术发展角度来看，
生成式人工智能正在逐渐改变数字内容生产方式。
随着扩散模型、多模态大模型以及三维生成技术的不断发展，
未来的三维内容生产将更加依赖自动化与智能化工具。
在这一背景下，
构建能够结合几何信息与参考材质信息的自动生成系统，
不仅能够提升三维资产生成质量，
还能够探索更加高效的三维生成技术路径。
通过将相关算法能力与 Web 平台结合，
还可以将复杂的生成技术转化为直观易用的工具，
进一步降低三维内容创作的技术门槛。

因此，本课题不仅在技术实现层面具有一定的研究价值，
同时在实际应用与产业发展方面也具有积极意义。
通过构建一个面向三维模型材质迁移任务的自动化系统，
本研究有望为三维内容生产流程提供更加高效、便捷的解决方案，
并为未来 AI 辅助三维设计工具的发展提供参考。